\documentclass[a4paper]{article}
\usepackage[pdftex]{graphicx}
\usepackage[utf8]{inputenc}
\usepackage{enumerate}
\usepackage{siunitx}
\sisetup{locale=DE} 
\usepackage{href-ul}
\hypersetup{
	colorlinks=true,
	linkcolor=blue,
	urlcolor=blue}
\usepackage{icomma}
\usepackage{amssymb}
\hyphenation{Ver-ein-fach-ung}
\usepackage{geometry}
\geometry{a4paper, top=15mm, left=15mm, right=15mm, bottom=15mm,
	headsep=10mm, footskip=12mm}
%Source: img_0146

% Start the document
\begin{document}
%Titel
{\bf Stegreifaufgabe aus der Physik }\\
\begin{enumerate}[1.]
%Aufgabe 1
{\bf \item Wie lautet der Impulserhaltungssatz?}
%Aufgabe 2
{\bf \item Ein Kleinbus der Masse $\SI{1,7}{\tonne}$ stößt frontal mit der Geschwindigkeit  $\SI{54}{\kilo\meter\over \hour}$  mit einem Kleinwagen der Masse $\SI{650}{\kilogram}$ zusammen,} der sich vor dem Zusammenprall mit der Geschwindigkeit  $\SI{36}{\kilo\meter\over \hour}$  bewegte. Der Stoß soll zur Vereinfachung als vollkommen unelastisch angenommen werden.

\begin{enumerate}[a)]
\item Berechnen Sie die Geschwindigkeit der beiden Wagen nach dem Zusammenstoß und geben Sie an, in welche Richtung sich die beiden Wagen bewegen. [falsches Ersatzergebnis: $ u= \SI{27}{\kilo\meter\over \hour}$]
\item Wie viel Prozent der vor dem Unfall vorhandenen kinetischen Energie wird durch den Stoß in 
Verformungsarbeit und innere Energie umgewandelt? 
\end{enumerate}

{\bf \item Bei der V2-Rakete strömten die Verbrennungsgase mit $\SI{2,0}{\kilo\meter\over \second}$ aus.} Damit entwickelte die Rakete eine Schubkraft von $\SI{2,0e5}{\newton}$.

\begin{enumerate}[a)]
\item Wie viel Treibstoff musste pro Sekunde verbrannt werden?
\item Die Masse der vollgetankten Rakete betrug $\SI{10}{\tonne}$. Mit welcher Beschleunigung hob die Rakete in lotrechter Richtung vom Startplatz ab?
\end{enumerate} 

{\bf \item Rückstoß: Beim $\alpha$-Zerfall eines Uranatoms (m = 238 u) wird ein $\alpha$-Teilchen (m = 4 u) aus dem Atomkern ausgestoßen.}
	Dabei erreicht es eine Geschwindigkeit von $\SI{1,4e7}{\meter\over \second}$
	\begin{enumerate}[a)]
		\item	Wie groß ist die Geschwindigkeit, mit der der Restkern zurückgestoßen wird?
		\item Wie ist das Verhältnis beider kinetischer Energien?
	\end{enumerate}
 
	
{\bf \item Deep Impact: Der Asteroid, der im Nördlinger Ries vor 15 Millionen Jahren einschlug, hatte eine Masse von ca. $\SI{7,0e13}{\kilogram}$ und eine Geschwindigkeit von
$\SI{20}{\kilo\meter\over \second}$ relativ zur Erde.} Berechne den Impuls, den er auf die Erde übertrug und entscheide, ob die Erde durch einen solchen Asteroiden aus der Bahn geworfen werden könnte (Masse Erde: $\SI{5,97e24}{\kilogram}$).  

\noindent
\begin{minipage}{0.5\textwidth}
{\bf \item Atwoodsche Fallmaschine}	
Die Massen $m_1=\SI{0,50}{\kilogram}$ und $m_2=\SI{0,40}{\kilogram}$ sind über eine Schnur und über eine Rolle (beide masselos) verbunden. Der Körper $m_2$ wird aus seiner Position P1 losgelassen. Die Schnur hat eine Schlaufe und wird erst gespannt, wenn die Masse $m_2$ $\SI{0,30}{\meter}$ tief frei gefallen ist. 
\begin{enumerate}[a)]
	\item Berechne die Geschwindigkeit, mit der die Masse $m_1$ dann angehoben wird.
	\item Welche Höhe $h_1$ erreicht die Masse $m_1$ ?
\end{enumerate}
\end{minipage}
\hfill
\begin{minipage}{0.5\textwidth}
	\includegraphics[width=7 cm]{atwo136}
\end{minipage}

 

\end{enumerate} 

\fbox{
	\begin{minipage}{0.5\textwidth}
		Zur Lösung bitte \href{https://www.okuyakl.de/phys/p10impL136/ll136.pdf}{hier klicken} oder den QR-Code scannen.\\
	Weitere Arbeitsblätter gibt es unter 
	
	\href{https://www.okuyakl.de}{www.okuyakl.de}
	\end{minipage}
	\hfill
	\begin{minipage}{0.4\textwidth}
		\includegraphics[width=1.5 cm]{../../viecher/zwe03}
		\includegraphics[width=3 cm]{qr136}
		\includegraphics[width=2 cm]{../../viecher/afanticon1}
		
	\end{minipage}}

\end{document}